\textbf{Hypotheses (fixed before the runs).} H1: CE memorises (mem\_rate $>0.9$) and loses more than 10 points of test accuracy at 40\% and 60\% noise relative to 0\% noise. H2: label smoothing ($\epsilon{=}0.1$) gains less than 3 points over CE at 40\% noise. H3: mixup ($\alpha{=}1$) lowers mem\_rate and raises test accuracy over CE at 20--60\% noise. H4: small-loss selection with the true noise rate is the most accurate system at 40\% and 60\% noise, with the most negative mem\_gap, but is worse than CE at 0\% noise. H5: small-loss selection depends on the assumed forget rate: under-estimating the noise rate loses most of the gain.

\textbf{Data.} scikit-learn digits (1,797 images, $8{\times}8$, $K{=}10$), pixels divided by 16. One fixed stratified split (\texttt{random\_state=1234}): 1,078 training and 719 test images; the split does not change with the seed. Symmetric noise is drawn on the training labels only, with a seed-dependent generator. Each noise rate is one task (digits\_noise0/20/40/60). There is no validation set and no early stopping: all hyperparameters were fixed a priori and the \emph{final-epoch} model is evaluated, so memorisation after long training is what is measured. The sweeps of Sec.~\ref{sec:ablations} are evaluated on the test split and are reported as sensitivity analyses; they were not used to choose the main settings.

\textbf{Model and training.} Conv(1$\to$16)--BN--ReLU, Conv(16$\to$32)--BN--ReLU--MaxPool, Conv(32$\to$64)--BN--ReLU--MaxPool, FC(256$\to$64)--ReLU--FC(64$\to$10), all convolutions $3{\times}3$ (about 41k parameters). SGD, momentum 0.9, learning rate 0.05 with cosine decay, weight decay $5\cdot10^{-4}$, batch size 64, 60 epochs, no data augmentation other than mixup where used. These settings are identical for all systems; the only per-system hyperparameters are $\epsilon{=}0.1$ (LS), $\alpha{=}1$ (mixup) and the schedule of Sec.~\ref{sec:method} (small-loss). No hyperparameter was tuned.

\textbf{Protocol.} 4 systems $\times$ 4 noise rates $\times$ 5 seeds (0--4) $=80$ main runs; the seed sets the noise draw, initialisation and batch order. The sensitivity runs use 40\% noise and 5 seeds each: 4 mixup $\alpha$ values, 4 LS $\epsilon$ values, 5 assumed forget rates and one no-warm-up variant (70 runs). Tables report mean $\pm$ standard deviation over seeds; $p$-values are paired two-sided $t$-tests over the 5 seeds (the seed fixes the noise draw, so the systems see the same corrupted labels). With $n{=}5$ these tests have little power and are indicative only; no multiple-comparison correction is applied.

\textbf{Compute.} CPU only, 2 threads, a shared machine; a run takes about 12\,s on an idle core and a median of 23\,s on the loaded machine used (the two slowest runs took \rhval{run/5cecfeef8b/runtime_s:0}\,s and \rhval{run/053a6b0376/runtime_s:0}\,s from contention); total about 80 CPU-minutes for 150 runs. All baselines are reimplemented here in one script with identical code paths for data, model and evaluation.
